Kurs:Algebraische Kurven (Osnabrück 2012)/Vorlesung 25/latex
Kurs:Algebraische Kurven (Osnabrück 2012)/Vorlesung 25/latex

\setcounter{section}{25}


  \zwischenueberschrift{Lösung in Potenzreihen für algebraische Kurven}


Es sei


\mavergleichskette

{\vergleichskette

{ F
}

{ \neq }{ 0
}

{  }{
}

{    }{
}

{   }{
}

}
{}{}{}
ein Polynom, das die ebene algebraische Kurve $C$ beschreibe, und sei


\mavergleichskette

{\vergleichskette

{ P
}

{ = }{ (0,0)
}

{ \in }{ C
}

{    }{
}

{   }{
}

}
{}{}{}
vorausgesetzt
\zusatzklammer {was keine Einschränkung ist, und durch Verschiebung immer erreicht werden kann} {} {.}
Wie kann man die Kurve im Nullpunkt mittels Potenzreihen beschreiben? Wann gibt es also einen durch nichtkonstante Potenzreihen
\mathkor {} {G} {und} {H} {}
mit konstantem Term $0$ definierten Ringhomomorphismus


\mathdisp {K[X,Y] \longrightarrow K[ \![T]\! ] \text{ mit } X \longmapsto G \text{ und } Y \longmapsto H} {  }


mit


\mavergleichskette

{\vergleichskette

{ F(G,H)
}

{ = }{ 0
}

{  }{
}

{    }{
}

{   }{
}

}
{}{}{}
\zusatzklammer {also einen Ringhomomorphismus
\maabb {} { K[X,Y]/(F) } { K [ \![ T ]\! ]
} {}} {} {?}
Es geht also um Lösungen der Gleichung


\mavergleichskettedisp

{\vergleichskette

{ F(X,Y)
}

{ =} { 0
}

{ } {
}

{ } {
}

{ } {
}

}
{}{}{}
in Potenzreihen, die das Verhalten der Kurve um die Punktlösung $(0,0)$ herum genauer beschreiben.


Der grundsätzliche Ansatz ist hier ein \stichwort {Potenzreihenansatz} {,} wie er beispielsweise auch in der Theorie der Differentialgleichungen verwendet wird. Man setzt


\mathdisp {G = {\sum }_{  k =0 }^{ \infty }  a _{  k  }  T  ^{  k  } \text{ und } H = {\sum }_{  \ell =0 }^{ \infty }  b _{  \ell  }  T  ^{  \ell  }} {  }


mit zunächst unbestimmten Koeffizienten $a_k$ und $b_\ell$ an. Das direkte Einsetzen in die beschreibende Gleichung


\mavergleichskette

{\vergleichskette

{ F
}

{ = }{ 0
}

{  }{
}

{    }{
}

{   }{
}

}
{}{}{}
und Ausmultiplizieren ergibt dann einen prinzipiell unendlichen Ausdruck. Allerdings ist für jedes $T^k$ der zugehörige Ausdruck für den Koeffizienten nur durch endlich viele Daten bestimmt, und zwar sind dafür nur die Koeffizenten von

\mathl{F,G}{} und $H$ unterhalb des Grades $k$ relevant. Da


\mavergleichskette

{\vergleichskette

{ F(G,H)
}

{ = }{ 0
}

{  }{
}

{    }{
}

{   }{
}

}
{}{}{}
sein soll, müssen die Koeffizienten von

\mathl{F,G,H}{} so sein, dass sich als Koeffizient zu $T^k$ stets $0$ ergibt.


Man sucht dann nach Bedingungen, wann es dafür Lösungen gibt, wie sie aussehen und ob sie eindeutig sind. Die Bedingung


\mavergleichskette

{\vergleichskette

{ a_0
}

{ = }{ b_0
}

{ = }{ 0
}

{    }{
}

{   }{
}

}
{}{}{}
ist dabei eine \stichwort {Anfangsbedingung} {,} die wiedergibt, dass die Potenzreihenlösung durch den Nullpunkt gehen soll.


Es ergibt sich schnell eine Bedingung an die linearen Terme der Potenzreihen
\zusatzklammer {also an $a_1$ und $b_1$} {} {,}
die man als eine weitere Rechtfertigung dafür ansehen kann, dass wir die Linearfaktoren des niedrigsten homogenen Bestandteiles $F_m$ in der homogenen Zerlegung von $F$ als Tangentengleichungen interpretiert haben.


\inputfaktbeweis

{Ebene algebraische Kurven/Potenzreihenlösung für Punkt/Linearer Term liegt auf Tangente/Fakt}

{Lemma}

{}

{


\faktsituation {}

\faktvoraussetzung {Es sei $K$ ein
\definitionsverweis {algebraisch abgeschlossener Körper}{}{}
und sei


\mavergleichskette

{\vergleichskette

{ F
}

{ \in }{ K[X,Y]
}

{  }{
}

{    }{
}

{   }{
}

}
{}{}{}
ein Polynom mit
\definitionsverweis {homogener Zerlegung}{}{}


\mavergleichskette

{\vergleichskette

{F
}

{ = }{ F_m + \cdots +  F_d
}

{  }{
}

{    }{
}

{   }{
}

}
{}{}{}
mit


\mavergleichskette

{\vergleichskette

{ d
}

{ \geq }{m
}

{ \geq }{ 1
}

{    }{
}

{   }{
}

}
{}{}{}
und


\mavergleichskette

{\vergleichskette

{ F_m
}

{ \neq }{ 0
}

{  }{
}

{    }{
}

{   }{
}

}
{}{}{.}
Es sei


\mavergleichskettedisp

{\vergleichskette

{ F_m
}

{ =} { \prod_{\lambda = 1}^{m} ( u_\lambda X + v_\lambda Y )
}

{ } {
}

{ } {
}

{ } {
}

}
{}{}{}
die Faktorzerlegung in Linearfaktoren
\zusatzklammer {diese Linearfaktoren definieren also die Tangenten von


\mavergleichskettek

{\vergleichskettek

{ C
}

{ = }{ V(F)
}

{  }{
}

{    }{
}

{   }{
}

}
{}{}{}
an


\mavergleichskettek

{\vergleichskettek

{ P
}

{ = }{ (0,0)
}

{  }{
}

{    }{
}

{   }{
}

}
{}{}{}} {} {.}
Es seien


\mathdisp {G=\sum_{n=0}^\infty a_n T^n \text{ und } H=\sum_{\ell=0}^\infty b_\ell T^\ell \in K [ \![ T]\! ]} {  }


Potenzreihen, die eine Lösung der Kurvengleichung


\mavergleichskette

{\vergleichskette

{F
}

{ = }{ 0
}

{  }{
}

{    }{
}

{   }{
}

}
{}{}{}
durch den Nullpunkt beschreiben
\zusatzklammer {d.h.


\mavergleichskettek

{\vergleichskettek

{ a_0
}

{ = }{ b_0
}

{ = }{ 0
}

{    }{
}

{   }{}

}
{}{}{}} {} {.}}

\faktfolgerung {Dann ist


\mavergleichskette

{\vergleichskette

{ u_\lambda a_1+ v_\lambda b_1
}

{ = }{ 0
}

{  }{
}

{    }{
}

{   }{
}

}
{}{}{}
für ein $\lambda$, d.h. der lineare Term der Potenzreihen ist durch eine der Tangenten vorgegeben.}

\faktzusatz {}

\faktzusatz {}


}

{


Wir setzen


\mathdisp {G=a_1T+a_2T^2+ \ldots \text{   und } H=b_1T+b_2T^2+ \ldots} {  }


in $F$ ein. In einem homogenen Bestandteil $F_k$, der ja eine Summe von Ausdrücken der Form \mathkon  { c_{ij}X^{i}Y^{j}  }  {  mit  }  { i+j=k  }{ } ist, kann man sofort $T^k$ ausklammern, und zwar ergibt sich ein Ausdruck der Form


\mavergleichskettedisphandlinks

{\vergleichskettedisphandlinks

{ F_k(G,H)
}

{ =} { \left(  \sum_{i+j = k} c_{ij} a_1^{i} b_1^{j}  \right)  T^k + \left(  \sum_{i+j = k} c_{ij} \left(  i a_1^{i-1}a_2b_1^{j} + j a_1^{i}b_1^{j-1}b_2  \right)  \right)   T^{k+1} + \ldots
}

{ } {
}

{ } {
}

{ } {
}

}
{}{}{.}
In den Koeffizient von $T^k$ gehen also

\mathl{a_1,b_1}{} in einer übersichtlichen Form über $F_k$ ein, aber auch kompliziertere Terme, die von \mathind { F_\ell  } {  \ell <k  }{,} herrühren. Auf $F_m$ angewandt, wo ja keine kleineren homogenen Komponenten mitberücksichtigt werden müssen, heißt dies, dass


\mavergleichskettedisp

{\vergleichskette

{ \sum _{i+j = m} c_{ij} a_1^{i} b_1^{j}
}

{ =} { 0
}

{ } {
}

{ } {
}

{ } {
}

}
{}{}{}
die entscheidende Gleichung für $a_1$ und $b_1$ ist. Dies ist aber nichts anderes als die Bedingung


\mavergleichskettedisp

{\vergleichskette

{ F_m(a_1,b_1)
}

{ =} { 0
}

{ } {
}

{ } {
}

{ } {
}

}
{}{}{.}
Da $F_m$ ein Produkt von Linearfaktoren ist, muss

\mathl{\left( a_1   , \,  b_1                   \right)}{} einen der Linearfaktoren annullieren, was die behauptete Aussage ist.


}


Man beachte, dass im vorstehenden Lemma die Möglichkeit


\mavergleichskette

{\vergleichskette

{ a_1
}

{ = }{ b_1
}

{ = }{ 0
}

{    }{
}

{   }{
}

}
{}{}{}
nicht ausgeschlossen ist. In der Tat gibt es nur unter zusätzlichen Bedingungen eine Realisierung einer Kurve mittels Potenzreihen entlang einer vorgegebenen Tangente, siehe
Satz 25.2
und die Beispiele weiter unten.


Der Rechenaufwand zur Bestimmung einer Potenzreihenlösung lässt sich wesentlich verringern, wenn man sich auf \anfuehrung{Graphenlösungen}{} beschränkt, wo die eine Potenzreihe einfach ein lineares Polynom ist
\zusatzklammer {und zwar eines, das durch eine Tangente gegeben ist} {} {,}
und die zweite eine zu bestimmende Potenzreihe. Das ist häufig keine wesentliche Einschränkung, wie aus
Lemma 24.11
folgt. Mit diesem Lemma können wir nämlich die Potenzreihen


\mavergleichskette

{\vergleichskette

{ G,H
}

{ \in }{ K[ \![T]\! ]
}

{  }{
}

{    }{
}

{   }{
}

}
{}{}{}
einfach transformieren, wenn nicht beide linearen Terme verschwinden. Sei hierzu


\mathbed {G = a_1T+ \ldots} {}

 {a_1 \neq 0} {}

 {} {} {} {}
angenommen. Mit einer geeigneten Potenzreihe

\mathl{U(T)}{} ist


\mavergleichskette

{\vergleichskette

{ U(G(T))
}

{ = }{ T
}

{  }{
}

{    }{
}

{   }{
}

}
{}{}{}
und


\mavergleichskette

{\vergleichskette

{ U(H(T))
}

{ = }{  \tilde{H}(T)
}

{  }{
}

{    }{
}

{   }{
}

}
{}{}{.}
Man schaltet also einen Potenzreihenautomorphismus dahinter, damit die Hintereinanderschaltung


\mathdisp {K[X,Y] \stackrel{X \mapsto G, Y  \mapsto H}{\longrightarrow} K[ \![T]\! ] \stackrel{T \mapsto U(T) }{ \longrightarrow} K[ \![T]\! ]} {  }


die besonders einfache Gestalt \mathkon  { X \mapsto T  }  {  ,  }  { Y \mapsto \tilde{H}   }{ } bekommt. Dies bedeutet, dass man die Kurve als Graph zu einer formalen Funktion in einer Variablen realisieren möchte.


\inputfaktbeweis

{Ebene algebraische Kurven/Tangenten mit Kontaktordnung eins/Formal-analytische Realisierung als Graph/Fakt}

{Satz}

{}

{


\faktsituation {}

\faktvoraussetzung {Es sei $K$ ein Körper und sei


\mavergleichskette

{\vergleichskette

{ F
}

{ \in }{ K[X, Y]
}

{  }{
}

{    }{
}

{   }{
}

}
{}{}{}
ein Polynom $\neq 0$ mit


\mavergleichskette

{\vergleichskette

{ (0,0)
}

{ \in }{ C
}

{ = }{ V(F)
}

{    }{
}

{   }{}

}
{}{}{}
und sei


\mavergleichskette

{\vergleichskette

{ F
}

{ = }{ F_d + \cdots + F_m
}

{  }{
}

{    }{
}

{   }{
}

}
{}{}{}
die
\definitionsverweis {homogene Zerlegung}{}{}
von $F$ mit


\mavergleichskette

{\vergleichskette

{d
}

{ \geq }{m
}

{  }{
}

{    }{
}

{   }{
}

}
{}{}{}
und mit


\mavergleichskette

{\vergleichskette

{F_m
}

{ \neq }{ 0
}

{  }{
}

{    }{
}

{   }{
}

}
{}{}{.}
Es sei

\mathl{uX+vY}{} ein einfacher Linearfaktor von $F_m$
\zusatzklammer {also ein lineares Polynom, das eine Tangente mit Multiplizität $1$ definiert} {} {.}}

\faktfolgerung {Dann gibt es Potenzreihen


\mathdisp {G= \sum_{n=0}^\infty a_n T^n,\,  H = \sum_{\ell=0}^\infty b_\ell T^\ell \in K [ \![ T]\! ]} {  }


mit


\mavergleichskette

{\vergleichskette

{ F(G,H)
}

{ = }{ 0
}

{  }{
}

{    }{
}

{   }{
}

}
{}{}{}
und mit konstantem Term


\mavergleichskette

{\vergleichskette

{ a_0,b_0
}

{ = }{ 0
}

{  }{
}

{    }{
}

{   }{
}

}
{}{}{}
und mit


\mavergleichskette

{\vergleichskette

{ a_1u+b_1v
}

{ = }{ 0
}

{  }{
}

{    }{
}

{   }{
}

}
{}{}{.}}

\faktzusatz {Dabei kann eine der Potenzreihen als ein lineares Polynom gewählt werden.}

\faktzusatz {}


}

{


Durch eine lineare Variablentransformation können wir annehmen, dass


\mavergleichskette

{\vergleichskette

{ uX+vY
}

{ = }{ Y
}

{  }{
}

{    }{
}

{   }{
}

}
{}{}{}
ist. Wir zeigen, dass es dann eine Potenzreihenlösung mit


\mavergleichskette

{\vergleichskette

{G
}

{ = }{ T
}

{  }{
}

{    }{
}

{   }{
}

}
{}{}{}
und mit
\zusatzklammer {zu konstruierendem} {} {}


\mavergleichskette

{\vergleichskette

{ H
}

{ = }{b_2T^2+b_3T^3+ \ldots
}

{  }{
}

{    }{
}

{   }{
}

}
{}{}{}
gibt. Wegen


\mavergleichskette

{\vergleichskette

{a_1
}

{ = }{ 1
}

{  }{
}

{    }{
}

{   }{
}

}
{}{}{}
und


\mavergleichskette

{\vergleichskette

{b_1
}

{ = }{ 0
}

{  }{
}

{    }{
}

{   }{
}

}
{}{}{}
erfüllt das die angegebene lineare
\zusatzklammer {Tangenten} {-} {}Bedingung.


Es sei


\mavergleichskette

{\vergleichskette

{ F
}

{ = }{ \sum_{i,j} c_{ij} X^{i} Y^{j}
}

{  }{
}

{    }{
}

{   }{
}

}
{}{}{.}
Es ist


\mavergleichskette

{\vergleichskette

{ c_{m,0}
}

{ = }{ 0
}

{  }{
}

{    }{
}

{   }{
}

}
{}{}{,}
da andernfalls $Y$ kein Linearfaktor von $F_m$ sein könnte, und es ist


\mavergleichskette

{\vergleichskette

{ c_{m-1,1}
}

{ \neq }{ 0
}

{  }{
}

{    }{
}

{   }{
}

}
{}{}{,}
da sonst $Y$ ein Linearfaktor mit einer Multiplizität $\geq 2$ wäre.


Wir zeigen, dass es bei diesen Anfangsdaten eine eindeutig bestimmte Potenzreihe


\mavergleichskette

{\vergleichskette

{H
}

{ = }{b_2T^2+b_3T^3+ \ldots
}

{  }{
}

{    }{
}

{   }{
}

}
{}{}{}
gibt. Einsetzen von $G$ und $H$ in $F$ ergibt für jedes $k$ eine Bedingung, da der resultierende Koeffizient zu $T^k$ gleich $0$ sein muss. Der $k$-te Koeffizient ist eine Summe von Ausdrücken der Form


\mathdisp {c_{ij} b_{\ell_1 } \cdots b_{\ell_j } \text{ mit } i+ \sum_{\rho=1}^{j}  \ell_\rho = k} {  }


\zusatzklammer {diese Ausdrücke kommen mehrfach mit einem gewissen Multinomialkoeffizienten vor} {} {.}
Da


\mavergleichskette

{\vergleichskette

{ \ell_\rho
}

{ \geq }{ 2
}

{  }{
}

{    }{
}

{   }{
}

}
{}{}{}
ist, kommt für


\mavergleichskette

{\vergleichskette

{k
}

{ < }{ m + \ell -1
}

{  }{
}

{    }{
}

{   }{
}

}
{}{}{}
der Term $b_\ell$ nicht vor. Der Term $b_\ell$ kommt erstmals im


\mavergleichskette

{\vergleichskette

{ k
}

{ = }{ (m+ \ell -1)
}

{  }{
}

{    }{
}

{   }{
}

}
{}{}{-}ten Koeffizienten vor, und zwar in der einzigen Weise


\mathdisp {c_{m-1,1} b_\ell} { . }


Ansonsten kommen in diesem Koeffizienten nur die $c_{ij}$ und $b_r$ mit


\mavergleichskette

{\vergleichskette

{r
}

{ < }{ \ell
}

{  }{
}

{    }{
}

{   }{
}

}
{}{}{}
vor. Da nach Voraussetzung


\mavergleichskette

{\vergleichskette

{ c_{m-1,1}
}

{ \neq }{ 0
}

{  }{
}

{    }{
}

{   }{
}

}
{}{}{}
ist, ist dadurch der Wert von $b_\ell$ eindeutig festgelegt. Die Koeffizienten $b_\ell$ werden also induktiv konstruiert, wobei die Werte jeweils eindeutig durch die Bedingung an die Koeffizienten festgelegt sind.


}


\inputbeispiel{ }

{


Wir betrachten die ebene affine Kurve vom Grad drei, die durch die Gleichung


\mavergleichskette

{\vergleichskette

{ F
}

{ = }{ X^3+XY+Y
}

{ = }{ 0
}

{    }{
}

{   }{
}

}
{}{}{}
gegeben ist. Die partiellen Ableitungen sind


\mathdisp {\frac{\partial F}{\partial X} =3 X^2 + Y \text{ und }  \frac{\partial F}{\partial Y} = X+1} { . }


Die zweite Ableitung ist nur bei


\mavergleichskette

{\vergleichskette

{ X
}

{ = }{ -1
}

{  }{
}

{    }{
}

{   }{
}

}
{}{}{}
gleich $0$, dort hat aber $F$ den Wert $-1$, d.h. die Kurve ist glatt. Im Nullpunkt haben die partiellen Ableitungen den Wert

\mathl{(0,1)}{.} Die zugehörige Tangente ist also die $X$-Achse, was dazu passt, dass der lineare Term der Kurvengleichung $Y$ ist.


Wir berechnen die Potenzreihe


\mavergleichskette

{\vergleichskette

{ Y
}

{ = }{ H(T)
}

{ = }{ {\sum }_{  \ell=0 }^{ \infty }  b _{  \ell }  T ^{  \ell }
}

{    }{
}

{   }{
}

}
{}{}{,}
die die Kurve im Nullpunkt als Graphen beschreibt
\zusatzklammer {es ist


\mavergleichskettek

{\vergleichskettek

{ X
}

{ = }{ T
}

{  }{
}

{    }{
}

{   }{
}

}
{}{}{}} {} {.}
Die Anfangsbedingungen sind


\mavergleichskette

{\vergleichskette

{b_0
}

{ = }{b_1
}

{ = }{ 0
}

{    }{
}

{   }{
}

}
{}{}{.}
Für die folgenden Koeffizienten von $H$ müssen wir aus der Gleichung


\mavergleichskettedisp

{\vergleichskette

{ F(T,H)
}

{ =} { T^3+TH+H
}

{ =} { 0
}

{ } {
}

{ } {
}

}
{}{}{,}
also


\mavergleichskettedisp

{\vergleichskette

{ T^3 + T( b_2T^2+b_3T^3 + \ldots ) + ( b_2T^2+b_3T^3 + \ldots )
}

{ =} { 0
}

{ } {
}

{ } {
}

{ } {
}

}
{}{}{}
über die Koeffizienten der $T^k$ die Bedingungen an $b_\ell$ herauslesen.


$b_2$. Der zweite Koeffizient
\zusatzklammer {der Gleichung} {} {}
liefert sofort


\mavergleichskette

{\vergleichskette

{ b_2
}

{ = }{ 0
}

{  }{
}

{    }{
}

{   }{
}

}
{}{}{.}


$b_3$. Der dritte Koeffizient liefert die Bedingung


\mavergleichskette

{\vergleichskette

{ 1+b_3
}

{ = }{ 0
}

{  }{
}

{    }{
}

{   }{
}

}
{}{}{,}
woraus


\mavergleichskette

{\vergleichskette

{b_3
}

{ = }{ -1
}

{  }{
}

{    }{
}

{   }{
}

}
{}{}{}
folgt.


Die folgenden Koeffizienten liefern die Bedingung


\mavergleichskette

{\vergleichskette

{ b_{\ell-1} + b_{\ell}
}

{ = }{ 0
}

{  }{
}

{    }{
}

{   }{
}

}
{}{}{,}
sodass also die folgenden $b_\ell$ abwechselnd $1$ und $-1$ sind. Man hat also eine einfache Rekursionsformel und es ist


\mavergleichskettedisp

{\vergleichskette

{ H
}

{ =} { -T^3+T^4-T^5+T^6-T^7 + \ldots
}

{ } {
}

{ } {
}

{ } {
}

}
{}{}{.}
Die Umformung der Kurvengleichung in


\mavergleichskettedisp

{\vergleichskette

{ Y
}

{ =} { \frac{-X^3}{1+X}
}

{ } {
}

{ } {
}

{ } {
}

}
{}{}{}
zeigt, dass hier der Graph einer rationalen Funktion
\zusatzklammer {mit einem Pol bei


\mavergleichskettek

{\vergleichskettek

{ X
}

{ = }{ -1
}

{  }{
}

{    }{
}

{   }{
}

}
{}{}{}} {} {}
vorliegt. Die angegebene Potenzreihe beschreibt also den Graphen einer rationalen Funktion als Graphen einer formal-analytischen Funktion.


}


\inputbeispiel{}

{


Wir betrachten
das Kartesische Blatt,
das durch


\mavergleichskette

{\vergleichskette

{ X^3+Y^3-3XY
}

{ = }{ 0
}

{  }{
}

{    }{
}

{   }{
}

}
{}{}{}
gegeben ist, im Nullpunkt und bezüglich der durch


\mavergleichskette

{\vergleichskette

{Y
}

{ = }{ 0
}

{  }{
}

{    }{
}

{   }{
}

}
{}{}{}
gegebenen Tangente und wollen die Potenzreihe bestimmen, mit der sich der \anfuehrung{Zweig}{} der Kurve, der diese Tangente definiert, als Graph beschreiben lässt. Wir setzen also


\mavergleichskette

{\vergleichskette

{X
}

{ = }{ T
}

{  }{
}

{    }{
}

{   }{
}

}
{}{}{}
und


\mavergleichskette

{\vergleichskette

{ H
}

{ = }{ b_2T^2+b_3T^3+b_4T^4+ \ldots
}

{  }{
}

{    }{
}

{   }{
}

}
{}{}{}
an und haben diese Koeffizienten zu bestimmen
\zusatzklammer {die Charakteristik von $K$ sei nicht $3$} {} {.}
Die Koeffizienten $b_\ell$ sind durch die Bedingung


\mavergleichskettealign

{\vergleichskettealign

{ 0
}

{ =} { T^3+ H^3-3TH
}

{ =} { T^3+(b_2T^2+b_3T^3 + \ldots )^3 -3T(b_2T^2+b_3T^3 + \ldots )
}

{ } {
}

{ } {
}

}
{}
{}{}
festgelegt. Das Einsetzen bzw. Ausmultiplizieren dieser Potenzreihe liefert zum ersten Mal für


\mavergleichskette

{\vergleichskette

{k
}

{ = }{ 3
}

{  }{
}

{    }{
}

{   }{
}

}
{}{}{}
eine Bedingung. Der Summand $X^3$
\zusatzklammer {bzw. $T^3$} {} {}
muss überhaupt nur einmal berücksichtigt werden, nämlich für


\mavergleichskette

{\vergleichskette

{k
}

{ = }{ 3
}

{  }{
}

{    }{
}

{   }{
}

}
{}{}{.}
Der Summand $Y^3$ muss erst ab


\mavergleichskette

{\vergleichskette

{k
}

{ \geq }{ 6
}

{  }{
}

{    }{
}

{   }{
}

}
{}{}{}
berücksichtigt werden, da ja


\mavergleichskette

{\vergleichskette

{Y
}

{ = }{ H
}

{  }{
}

{    }{
}

{   }{
}

}
{}{}{}
ein Vielfaches von $T^2$ ist. Der Summand $XY$ muss ab


\mavergleichskette

{\vergleichskette

{k
}

{ = }{ 3
}

{  }{
}

{    }{
}

{   }{
}

}
{}{}{}
berücksichtigt werden.


$b_2$. Hier hat man die Bedingung


\mavergleichskettedisp

{\vergleichskette

{ 1 -3 b_2
}

{ =} { 0
}

{ } {
}

{ } {
}

{ } {
}

}
{}{}{,}
woraus sich


\mavergleichskette

{\vergleichskette

{ b_2
}

{ = }{ \frac{1}{3}
}

{  }{
}

{    }{
}

{   }{
}

}
{}{}{}
ergibt.


$b_3$. Dies taucht erstmals in der Bedingung für den vierten Koeffizienten auf. Dort steht es aber isoliert, sodass


\mavergleichskette

{\vergleichskette

{b_3
}

{ = }{ 0
}

{  }{
}

{    }{
}

{   }{
}

}
{}{}{}
sein muss.


$b_4$. Aus dem gleichen Grund ist


\mavergleichskette

{\vergleichskette

{b_4
}

{ = }{ 0
}

{  }{
}

{    }{
}

{   }{
}

}
{}{}{.}


$b_5$. Hierfür ist der sechste Koeffizient entscheidend, und dabei ist jetzt auch $Y^3$ zu berücksichtigen. Es ergibt sich die Bedingung


\mavergleichskettedisp

{\vergleichskette

{ b_2^3-3 b_5
}

{ =} { 0
}

{ } {
}

{ } {
}

{ } {
}

}
{}{}{,}
also


\mavergleichskette

{\vergleichskette

{b_5
}

{ = }{ \frac{1}{81}
}

{  }{
}

{    }{
}

{   }{
}

}
{}{}{.}


\mathl{b_6, b_7,b_8}{.} Der Summand


\mavergleichskettedisp

{\vergleichskette

{ Y^3
}

{ =} { \left(  b_2T^2 +b_5T^5+ \ldots  \right) \left(  b_2T^2 +b_5T^5+ \ldots  \right) \left(  b_2T^2 +b_5T^5+ \ldots  \right)
}

{ } {
}

{ } {
}

{ } {
}

}
{}{}{}
leistet erstmals wieder für den neunten Koeffizienten einen Beitrag, und zwar ist dieser

\mathl{3b_2^2b_5}{.} Dies bedeutet, dass $b_6$ und $b_7$ isoliert stehen und daher null sein müssen. Für $b_8$ ergibt sich die Bedingung


\mavergleichskettedisp

{\vergleichskette

{ 3b_2^2b_5 -3 b_8
}

{ =} { 0
}

{ } {
}

{ } {
}

{ } {
}

}
{}{}{}
und daher ist


\mavergleichskette

{\vergleichskette

{b_8
}

{ = }{\frac{1}{729}
}

{  }{
}

{    }{
}

{   }{
}

}
{}{}{.}


Der Anfang der Potenzreihe $H$, der den einen Kurvenzweig im Nullpunkt als einen Graphen beschreibt, ist also


\mavergleichskettedisp

{\vergleichskette

{ H
}

{ =} { \frac{1}{3} T^2 + \frac{1}{81} T^5 + \frac{1}{729}T^8 + \ldots
}

{ } {
}

{ } {
}

{ } {
}

}
{}{}{.}


}


\inputbeispiel{}

{


Wir betrachten die durch


\mavergleichskette

{\vergleichskette

{X^3-Y^2
}

{ = }{ 0
}

{  }{
}

{    }{
}

{   }{
}

}
{}{}{}
definierte
\definitionsverweis {Neilsche Parabel}{}{.}
Hier ist der Nullpunkt singulär, und es gibt nur eine Tangente, nämlich


\mavergleichskette

{\vergleichskette

{Y
}

{ = }{ 0
}

{  }{
}

{    }{
}

{   }{
}

}
{}{}{,}
diese hat aber die Multiplizität zwei, d.h.
Satz 25.2
ist hier nicht anwendbar. Wir werden zeigen, dass es überhaupt keine Potenzreihenlösung im Nullpunkt mit nicht verschwindendem linearen Term gibt.


Es seien dazu


\mavergleichskette

{\vergleichskette

{X
}

{ = }{ G
}

{ = }{ a_1T+a_2T^2 + \cdots
}

{    }{
}

{   }{
}

}
{}{}{}
und


\mavergleichskette

{\vergleichskette

{Y
}

{ = }{ H
}

{ = }{ b_1T+b_2T^2 + \cdots
}

{    }{
}

{   }{
}

}
{}{}{}
Potenzreihen, die die Kurvengleichung erfüllen. Wir setzen in die Kurvengleichung ein und erhalten für den zweiten Koeffizient die Bedingung


\mavergleichskette

{\vergleichskette

{ -b_1^2
}

{ = }{ 0
}

{  }{
}

{    }{
}

{   }{
}

}
{}{}{,}
woraus


\mavergleichskette

{\vergleichskette

{b_1
}

{ = }{ 0
}

{  }{
}

{    }{
}

{   }{
}

}
{}{}{}
folgt. Für den dritten Koeffizienten ergibt sich hingegen


\mavergleichskette

{\vergleichskette

{a_1^3
}

{ = }{ 0
}

{  }{
}

{    }{
}

{   }{
}

}
{}{}{,}
also wieder


\mavergleichskette

{\vergleichskette

{a_1
}

{ = }{ 0
}

{  }{
}

{    }{
}

{   }{
}

}
{}{}{.}


Dennoch gibt es Potenzreihenlösungen für die Neilsche Parabel durch den Nullpunkt. Hierzu kann man einfach die monomiale Lösung \mathkon  { G=T^2  }  {  und  }  { H=T^3  }{ } nehmen
\zusatzklammer {die ja sogar eine Bijektion zwischen der affinen Geraden und der Neilschen Parabel stiftet} {} {.}
Der lineare Term davon ist freilich gleich $0$.


}


\inputbemerkung

{}

{


Es sei


\mavergleichskette

{\vergleichskette

{ K
}

{ = }{ \R
}

{  }{
}

{    }{
}

{   }{
}

}
{}{}{}
oder $= {\mathbb C}$ und


\mavergleichskette

{\vergleichskette

{ F
}

{ \in }{ K[X,Y]
}

{  }{
}

{    }{
}

{   }{
}

}
{}{}{.}
Bei


\mavergleichskettedisp

{\vergleichskette

{ \left( \frac{\partial F}{\partial x}(P),\frac{\partial F}{\partial y} (P) \right)
}

{ \neq} { (0,0)
}

{ } {
}

{ } {
}

{ } {
}

}
{}{}{,}


wenn also $P$ ein
\definitionsverweis {regulärer Punkt}{}{}
der Funktion $F$ ist
\zusatzklammer {oder, äquivalent, ein
\definitionsverweis {glatter Punkt}{}{}
von


\mavergleichskettek

{\vergleichskettek

{ C
}

{ = }{ V(F-F(P))
}

{  }{
}

{    }{
}

{   }{
}

}
{}{}{}} {} {,}
so sichert der
Satz über implizite Funktionen,
dass sich die Kurve in einer
\zusatzklammer {metrischen} {} {}
Umgebung des Punktes als Graph einer differenzierbaren Funktion darstellen lässt.


}
